% Section draft: the elementary route to a linear threshold.
% To be \input after the dense-completion section once the constant is settled.

\section{An elementary route to the linear threshold}\label{sec:linearroute}

Theorem~\ref{thm:main} rests on the abelian reading of
\cite{MuyesserPokrovskiy2025} through Lemma~\ref{lem:dense}.  In this section
we record how far our reduction removes that dependence by elementary means.  The
outcome is that Lemma~\ref{lem:dense} reduces to a single statement about
permutations of an interval, which we prove in the two extreme cases and in a
family of composite cases, and which we have verified exactly for every order
up to $30$.  We state it as a conjecture; a proof would make
Theorem~\ref{thm:main} independent of \cite{MuyesserPokrovskiy2025} and would
give an explicit value of $C$.

\subsection{From the residual set to Heffter's difference problem}

After Lemma~\ref{lem:realization} the special support is $\pm I$ with
$I\subseteq[1,M]$ and $M=A(D+1)$, and the residual set is
$X=\pm([1,K]\setminus I)$ in $\Z_n$, $n=2K+1$.  In the branch of
Lemma~\ref{lem:dense} that needs the completion theorem, $X$ has to be
partitioned into signed zero-sum triples $\{\pm a,\pm b,\pm c\}$.  Writing a
triple by its three magnitudes $a<b<c$ in $[1,K]$, the two possibilities are
$a+b=c$ and $a+b+c=n$; the first is a pair $\{b,c\}$ at distance $a$ and the
second is the pair $\{b,-c\}$ at cyclic distance $a$ in $\Z_n$.  For
$I=\emptyset$ this is Heffter's first difference problem when $n\equiv1\pmod
6$, solved for every $n$ by Peltesohn~\cite{Peltesohn}, and the general case
asks for the same partition of $[1,K]\setminus I$, that is, for Heffter's
problem with the small differences $I$ forbidden.  The hardest case in our experiments is
$I=[1,M]$, and it is the one we analyse.

\subsection{Folded sum permutations}

Let $K-M=3t$.  In every solution our searches found, the $t$ smallest elements
$M+1,\dots,M+t$ can be taken as the distances $a$, and the remaining $2t$
elements split into a lower half $M+t+1,\dots,M+2t$ and an upper half
$M+2t+1,\dots,K$ with each pair $\{b,c\}$ taking $b$ from the lower and $c$
from the upper half.  Writing $a=M+1+i$, $b=M+t+1+j$, $c=M+2t+1+k$ with
$i,j,k\in[0,t-1]$ and $s=M+1-t$, the two triple types become
\[
  c-b=a\iff k=i+j+s,\qquad b+c=n-a\iff k=(2t-1)-(i+j+s).
\]
So a solution is a permutation $\pi$ of $[0,t-1]$, $j=\pi(i)$, such that the
numbers $u_i=i+\pi(i)+s$ lie in $[0,2t-1]$ and their folds
$\min(u_i,\,2t-1-u_i)$ form a permutation of $[0,t-1]$.  We call such a $\pi$
a \emph{folded sum permutation} of order $t$ and shift $s$, and write
$\mathrm{FSP}(t,s)$.  Equivalently, with $\tau_i=i+\pi(i)$, the $\tau_i$ are
distinct, lie in the window $W=[-s,2t-1-s]$, and no two of them sum to the
reflection constant $c=2t-1-2s$ of $W$; that is, the sum set of $\pi$ is a
transversal of the reflection of $W$.  Since $\sum_i\tau_i=t(t-1)$ for every
permutation, the mean of the transversal exceeds the centre of $W$ by
$s-\tfrac12$.

The range $3M+1\le K\le 7M+3$ of Lemma~\ref{lem:dense} corresponds to
$(1-t)/2\le s\le(t+1)/2$, and beyond $K\ge7M+3$ Simpson's theorem on Langford
sequences applies with no folding at all.

\begin{conjecture}\label{conj:fsp}
For every $t\ge5$ and every integer $s$ with $(1-t)/2\le s\le(t+1)/2$ there is
a folded sum permutation of order $t$ and shift $s$.
\end{conjecture}

The window is sharp: outside it no $\mathrm{FSP}(t,s)$ exists, because the
$s$ largest values of $W$ (for $s\ge1$) have no reflection partner in the
range of the $\tau_i$ and must all be attained, which is impossible once
$s>(t+1)/2$.  Our constraint programs verify the conjecture exactly, for every $3\le t\le30$ except $t=4$, where $s=-1$ and $s=2$
fail; unlike Simpson's theorem there are no congruence exceptions inside the
window, the folding removes them.

\subsection{What is proved}

\begin{proposition}\label{prop:fspcases}
Conjecture~\ref{conj:fsp} holds in the following cases.
\begin{enumerate}
\item $s\in\{0,1\}$: the identity is an $\mathrm{FSP}(t,s)$.
\item $t$ odd and $s\in\{(1-t)/2,(t+1)/2\}$: the rotation
  $\pi(i)=i+(t-1)/2 \bmod t$ is an $\mathrm{FSP}(t,s)$ at both endpoints.
\item (Composition.)  If $a$ is odd and $\sigma$ is an
  $\mathrm{FSP}(b,s')$, then
  \[
    \pi(au+w)=a\,\sigma(u)+\bigl((w+\tfrac{a-1}{2})\bmod a\bigr),\qquad
    u\in[0,b-1],\ w\in[0,a-1],
  \]
  is an $\mathrm{FSP}(ab,s)$ with $s=a(s'-1)+(a+1)/2$.
\end{enumerate}
\end{proposition}

\begin{proof}
(1) The sums $2i$ are distinct and even, the constant $c$ is odd, and
$0\le 2i\le 2t-2\le 2t-1-s$ exactly when $s\le1$; the lower bound $2i\ge-s$
needs $s\ge0$.

(2) For the rotation the sums $i+\pi(i)$ are $2i+(t-1)/2$ for
$i<(t+1)/2$ and $2i-(t+1)/2$ for $i\ge(t+1)/2$, which together form the
interval $[(t-1)/2,(3t-3)/2]$.  At $s=(1-t)/2$ this interval is the lower half
of $W$, at $s=(t+1)/2$ it is the upper half, and a half of $W$ contains no
reflected pair.

(3) On each block the inner rotation has sums $w+\rho(w)$ forming the interval
$I=[(a-1)/2,(3a-3)/2]$, which is carried to itself by $x\mapsto 2a-2-x$.  The
sums of $\pi$ are $a(u+\sigma(u))+I$.  Put $s=(a+1)/2+aq$.  The condition
$a(u+\sigma(u))+I\subseteq[-s,2ab-1-s]$ is $u+\sigma(u)\in[-(q+1),2b-q-2]$,
the window of $\mathrm{FSP}(b,q+1)$; and the reflection constant
$c=2ab-1-2s=a(2b-2q-1)-2$ sends $a\,y+x$ to $a(2b-2q-3-y)+(2a-2-x)$, so it
reflects $I$ onto itself and the block index $y$ by the constant
$2b-1-2(q+1)$ of $\mathrm{FSP}(b,q+1)$.  Therefore the sum set of $\pi$ is a
transversal exactly when the sum set of $\sigma$ is, and $s'=q+1$.
\end{proof}

Case (2) has two readings: at the lower endpoint it is a hooked Langford
sequence of defect $(t+1)/2$, at the upper endpoint the columns
$(i,\pi(i),c-i-\pi(i))$ form a $3\times t$ magic rectangle.  Case (3) covers
the residues $s\equiv(a+1)/2\pmod a$ for every odd divisor $a$ of $t$; it does
not reach prime $t$, and we have not found an additive counterpart.

\subsection{Consequences}

If Conjecture~\ref{conj:fsp} holds then the clean tail $[M+1,K]$ has a
maximum packing into difference triples for every $K\ge3M+1$ (the cases
$K-M\not\equiv0\pmod3$ reduce to it by omitting the one or two smallest
elements, and the finitely many cases $K\le3M+3$ are magic rectangles with one
or two empty cells, verified for $M\le30$).  Together with the proved branch of
Lemma~\ref{lem:dense} for $K\ge7M+3$, this gives Lemma~\ref{lem:dense} with
$\eta=1/(14A+O(1))$, hence Theorem~\ref{thm:main} with an explicit constant
$C$ and without any use of \cite{MuyesserPokrovskiy2025}.  The punctured
cases $I\subsetneq[1,M]$ are the same object with the extra small elements
as additional distances, and were found feasible for every random $I$ we
tried; we have not reduced them to Conjecture~\ref{conj:fsp} in general.
